\section{Algebraic degrees of the extremal constants}

\begin{theorem}\label{th:degree}
The field $K=\mathbb Q(\tp^2)$ has degree $48$ and equals
$\mathbb Q(\kappa)$. Each of $\tp$ and $C$ has degree $96$ over $\mathbb Q$.
\end{theorem}

\begin{proof}
Set $u=\tp^2$ and $b_j=30w_j$ for $1\le j\le6$, so
$(b_1,\ldots,b_6)=(30,15,20,15,24,10)$. Introduce two
parameters $s,t$ and put
\[
 B_j=\bigl(b_j(4j^2-x)+s\,\mathbf1_{j=1}
              +t b_j(x-j^2)\bigr)\sqrt{x-j^2}.
\]
Write $\Omega(x,s,t)$ for the product of $B_1+\sum_{j=2}^6\varepsilon_j B_j$
over the $32$ choices of signs. Changing any radical sign preserves the product.
For the first sign change, negate all $32$ factors and permute them. Applying
this argument with independent formal radicals shows that each radical has
even exponent. Thus $\Omega\in\Z[x,s,t]$ and its degree in $x$ is at most $48$.
Set $P(x)=\Omega(x,0,0)$, and let $H_s,H_t$ be the two parameter derivatives
at $s=t=0$. On the positive branch at $x=u$ the all-plus factor vanishes by
the critical-point equation for $F$. In particular $P(u)=0$.

The accompanying exact certificate gives $\deg(P\bmod 53)=48$ and shows
that this reduction is irreducible. Here are the details of the calculation
for independent reproduction. Evaluate the three polynomials $P,H_s,H_t$ at
$x=0,\ldots,48$ over $\mathbb F_{53}$ and interpolate. To compute a product
over signs, successively multiply by its conjugate under each of the last
five radical sign changes. In the basis indexed by subsets $I\subseteq
\{1,\ldots,6\}$, multiplication replaces the two subsets $I,J$ by
$I\mathbin\triangle J$ and contributes a factor
$\prod_{j\in I\cap J}(x-j^2)$. The product rule carries both derivatives.
In $\mathbb F_{53}[X]$, writing $\bar P$ for the reduction, the certificate
checks
\[
 X^{53^{48}}\equiv X\pmod{\bar P},\qquad
 \gcd(\bar P,X^{53^d}-X)=1\quad(d=16,24).
\]
Since $2,3$ are the prime divisors of $48$, these tests exclude every
proper possible irreducible-factor degree and prove irreducibility. Gauss's
lemma now gives $[K:\mathbb Q]=48$.

The root $u$ of $P$ is simple. Every radical is analytic near this real
point, so exactly one sign factor vanishes. Differentiating the product
and cancelling the nonvanishing factors gives
\[
 \frac{H_t(u)}{H_s(u)}
 =\frac{\sum_{j=1}^6 b_j(u-j^2)^{3/2}}{\sqrt{u-1}},
 \qquad
 \kappa=\frac{2025u^4 H_s(u)^2}{256(u-1)H_t(u)^2}.
\]
Both derivatives at $u$ are nonzero. This places $\kappa$ in $K$.
In $\mathbb F_{53}[X]/(\bar P)$, the denominator of the formula is
nonzero. The certificate verifies by cross multiplication that neither
the $16$th nor the $24$th iterate of Frobenius fixes the quotient.
It therefore has degree $48$ over $\mathbb F_{53}$. To justify reduction,
the leading coefficient of $P$ is a $53$-adic unit. After dividing by it,
$P$ is monic over $\Z_{(53)}$, so reduction of $\Z_{(53)}[u]$
gives the stated finite field. The map extends upon inverting the displayed
denominator. Any rational polynomial relation on $\kappa$ can be scaled to
a primitive integer relation, whose reduction is a nonzero polynomial
annihilating the finite-field quotient. Consequently $\kappa$ has degree
at least $48$, proving $\mathbb Q(\kappa)=K$.

Finally let $L$ be the leading coefficient of $P$. Since its degree is even,
the norms of $u$ and $u-1$ are $P(0)/L$ and $P(1)/L$. The certificate
gives their residues $50\pmod {53}$ and $42\pmod {59}$, respectively,
both quadratic nonresidues at primes where $L$ is a unit. Neither norm is
a rational square, hence neither $u$ nor $u-1$ is a square in $K$.
Now $\tp=\sqrt u$, and
$C=2H_t(u)\sqrt{u-1}/(45u^2 H_s(u))$. Since $\mathbb Q(C)$ contains
$\kappa$, it contains $K$. Both claimed quadratic extensions follow.
\end{proof}
